\chapter{Conclusion}
\label{Chapter7}
\lhead{Chapter 7. \emph{Conclusion}}

%─────────────────────────────────────────────────────────────────────────────
\section{Summary}
\label{sec:conclusion-summary}
%─────────────────────────────────────────────────────────────────────────────

CausalTrace defends a container cluster on two surfaces. The kernel
surface enforces physical invariants---a process whose standard input
points to a socket is a reverse shell regardless of how it got there. The
coupling surface observes inter-container behaviour as a cellular sheaf
whose Rayleigh quotient rises when the graph of service interactions
violates its calibrated structure. The two surfaces are independent; a
failed Tier-3 daemon never disables Tier-1 enforcement.

Four claims distinguish the work. The \emph{three-tier architecture with
an independence invariant} separates the hot path from the analytics
path. The \emph{compound enforcement gate with a four-level trust model}
replaces the binary kill/allow decision with a three-case taxonomy gated
by per-IP trust, and uses a TC classifier to drop only the attacker's
packets at the veth. The \emph{rank-$15$ CCA calibration pipeline}
produces a finite, well-conditioned coupling model on realistic windows
where the earlier $k{=}50$ iteration was rank-deficient. The
\emph{seed-reproducible evaluation harness} replays an identical
155-injection sequence against every detector under the same workload;
the resulting headline numbers are CausalTrace 99.4\,\% [96, 100],
Falco tuned 70.3\,\% [63, 77], Tetragon tuned 1.3\,\% [0, 5].

%─────────────────────────────────────────────────────────────────────────────
\section{Limitations}
\label{sec:conclusion-limitations}
%─────────────────────────────────────────────────────────────────────────────

Five limitations bound the results. \emph{Host root compromises the
defender}; CausalTrace is kernel-resident, and an attacker with host root
can unload the BPF programs. \emph{Results come from a single evaluation
host} on Linux~6.17, so cross-kernel generalisation remains unverified.
\emph{Probe~B is IPv4-only}; deployments with IPv6-first meshes need an
additional \texttt{tcp\_v6\_connect} hook. \emph{The trust model is
pessimistic for accept-thread-plus-thread-pool servers}, where per-TID
attribution falls back to the per-cgroup best-effort path. \emph{Baseline
tuning introduces a pro-CausalTrace bias}: Falco and Tetragon were tuned
for this attack set, and a production Falco operator would invest more
effort than our tuned ruleset reflects.

%─────────────────────────────────────────────────────────────────────────────
\section{Future work}
\label{sec:conclusion-future}
%─────────────────────────────────────────────────────────────────────────────

Each limitation flows to an extension.

\paragraph{Multi-host sheaf.} The cellular sheaf construction extends to
a bipartite graph with intra-host edges (the current case) and inter-host
edges (flow-level attributions between hosts). A distributed calibrator
that synchronises restriction maps per inter-host edge plus a gossip
protocol for burn-list synchronisation would lift CausalTrace from
single-node to cluster scale without new mathematics.

\paragraph{Incremental online calibration.} A rolling window that updates
the restriction maps under the same six-cycle pristine-streak rule the
guarded EMA already uses (Section~\ref{sec:ema}) would remove the
calibration barrier and let the system track workload drift without an
offline pass.

\paragraph{IPv6 and accept-thread attribution.} Both are engineering
extensions rather than research; the \texttt{tcp\_v6\_connect} hook is a
one-day patch, and a thread-pool-aware attribution path would require
the dispatcher to consult a per-thread-pool client-IP map populated by
the accept-thread itself.

\paragraph{Hardware offload for the TC classifier.} Moving the
\texttt{drop\_ip\_list} consultation into smart-NIC XDP offload would cut
per-packet cost from $\approx\!100$\,ns to sub-10\,ns and reduce CPU
load during a sustained attack.

\paragraph{Formal verification of the compound gate.} The three-case
taxonomy is stated informally in Section~\ref{sec:compound}. A formal
model in a proof assistant would let us prove invariants such as
``every Case-A trigger either burns the attacker's trust or issues a
SIGKILL'' and ``the gate never kills a calibrated client without a
Case-A invariant''. The gate is under 100 lines of C; verification is
tractable.

%─────────────────────────────────────────────────────────────────────────────
\section{Closing remarks}
\label{sec:conclusion-closing}
%─────────────────────────────────────────────────────────────────────────────

Container runtime security is a crowded field. What distinguishes
CausalTrace is the commitment to a two-surface formulation: strict
physical invariants enforced at the kernel boundary, plus a mathematical
model of inter-container coupling that detects the composite attacks no
single-container rule catches. The mathematics is classical---CCA,
Mahalanobis distance, cellular sheaves, count-min sketches. The
engineering is kernel-resident---a tail-call dispatcher, six specialist
handlers, TC direct-action classifiers, a compound gate with an explicit
three-case taxonomy.

The demonstration is that the two surfaces, composed carefully, are both
faster and more discriminating than the best-available production
tooling on the attack classes this field cares about. Whether the
system is \emph{deployable} is a question for operators and for the
extensions above. Whether it \emph{works} is answered in
Chapter~\ref{Chapter6}.
